% !TeX program = xelatex
% UTF-8; standalone XeLaTeX/XePersian, or input into a XePersian document.
% Faithful transcription; medical statements have not been revised.
% Google Drive: Notebook/seizure & epilepsy (six handwritten pages).
% pic01.jpg = 20261016_205412191.jpg
% pic02.jpg = 20261016_205420764.jpg
% pic03.jpg = 20261016_205425083.jpg
% pic04.jpg = 20261016_205430161.jpg
% pic05.jpg = 20261016_205434052.jpg
% pic06.jpg = 20261016_205438438.jpg
\ifdefined\settextfont
  \def\SeizureStandalone{0}
\else
  \def\SeizureStandalone{1}
  \documentclass[11pt,a4paper]{article}
  \usepackage[margin=2cm]{geometry}
  \usepackage{amsmath,amssymb,enumitem}
  \usepackage{xepersian}
  \IfFontExistsTF{Amiri}{\settextfont{Amiri}}{\settextfont{FreeSerif}}
  \setlatintextfont{FreeSerif}
  \setlist[itemize]{itemsep=2pt,topsep=3pt,parsep=0pt}
\fi
\providecommand{\en}[1]{\lr{#1}}
\newcommand{\szarrow}{\ensuremath{\leftarrow}}
\newcommand{\szunclear}[1]{\textbf{[خوانش نامطمئن: #1]}}
\newcommand{\szpage}[2]{%
  \clearpage\section*{۵ ـ تشنج و صرع — صفحهٔ #1}
  {\footnotesize\noindent\en{#2}\par
  \noindent\en{Subject:} ۵ ـ تشنج و صرع\hfill\en{Date: / /}\par
  \noindent\en{Sat.\quad Sun.\quad Mon.\quad Tue.\quad Thu.\quad Wed.\quad Fri.}\par}
  \medskip}
\ifnum\SeizureStandalone=1 \begin{document}\fi

\szpage{۱}{pic01.jpg — 20261016\_205412191.jpg}
\begin{center}\szunclear{بسمه تعالی}\end{center}
\subsection*{تشنج پارشیل}
\begin{itemize}
  \item پارشیل ساده \szarrow{} هوشیاری به طور کامل حفظ شده.
  \item پارشیل کمپلکس \szarrow{} با اختلال هوشیاری همراه است.
\end{itemize}
\subsection*{پارشیل ساده}
\begin{itemize}
  \item موجب اختلال حرکتی، حسی، روانی یا اتونوم می‌شوند.
  \item بدون اختلال هوشیاری.
  \item در \en{EEG} امواج غیرطبیعی در ناحیهٔ محدودی از کورتکس.
  \item اگر تشنج‌های کانونی از نقطهٔ محدودی آغاز شده مثل انگشتان، سپس به مناطق دیگر گسترش یابد \szarrow{} \en{Jacksonian Marsh}.
  \item سایر تظاهرات:
  \begin{itemize}
    \item تغییر حس سوماتیک: پارستزی.
    \item اختلالات بینایی: مشاهدهٔ نور، توهم بینایی.
    \item اختلالات تعادلی: احساس سقوط، سرگیجه.
    \item اختلال اتونوم: تپش قلب، سیخ شدن موها.
  \end{itemize}
  \item اگر منشأ حمله لوب فرونتال یا تمپورال باشد \szarrow{} اختلال حافظه حین حمله \szarrow{} آفازی بعد از حمله.
\end{itemize}
\subsection*{پارشیل کمپلکس}
\begin{itemize}
  \item منشأ: بخش تحتانی و داخلی لوب تمپورال / مناطق مجاور آن در لوب فرونتال.
  \item تظاهرات: تغییر گذرای هوشیاری (مهم‌ترین)؛ بیمار اتفاقات حین حمله را به خاطر ندارد ولی شروع حمله را ممکن است / توهم بویایی در شروع حمله (اورا).
  \item اورا چند ثانیه تا چند دقیقه طول می‌کشد (به طور کلی).
  \item اتوماتیسم: حرکات غیرارادی حین حمله یا بلافاصله پس از آن.
  \begin{itemize}
    \item بیمار آن را به خاطر نمی‌آورد.
    \item نگاه‌کردن، ضربه‌زدن آرام با دست‌ها، جویدن، قورت‌دادن، ملچ‌ملچ‌کردن لب‌ها، زیر لب حرف‌زدن.
  \end{itemize}
  \item معمولاً در حملات لوب تمپورال.
\end{itemize}

\szpage{۲}{pic02.jpg — 20261016\_205420764.jpg}
\begin{itemize}
  \item \en{EEG}: در بین حملات \en{N/L} یا امواج سوزنی و نوک‌تیز.
  \begin{itemize}
    \item در حین حمله هم ممکن است \en{N/L} باشد.
  \end{itemize}
\end{itemize}
\subsection*{تشنج پارشیل با انتشار ثانویه}
\begin{itemize}
  \item ممکن است حمله به هر دو نیمکره انتشار یافته و جنرالیزه شود \szarrow{} حمله تونیک ـ کلونیک است.
  \item معمولاً به دنبال یک حمله پارشیل ساده با منشأ لوب فرونتال.
  \item \en{Postictal Todd's Paralysis}:
  فلج تاد: فلج ایجادشده بعد از حملهٔ تشنج فوکال یا جنرالیزه.
  \begin{itemize}
    \item ضعف، فلج، گذرا، خودمحدودشونده.
  \end{itemize}
\end{itemize}
\subsection*{درمان تشنج‌های پارشیل}
\begin{itemize}
  \item به یک اندازه مؤثر هستند.
  \item کاربامازپین یا فنی‌توئین \szarrow{} تشنج‌های پارشیل و تشنج‌های پارشیلی که جنرالیزه می‌شوند.
  \item کاربامازپین \szarrow{} باید ۲–۳ بار در روز مصرف شود.
  \begin{itemize}
    \item رابطه بین دوز و سطح سرمی و سطح توکسیک: خطی \szarrow{} مزیت.
  \end{itemize}
  \item فنی‌توئین کینتیک اشباعی دارد.
  \item عوارض: لکوپنی، آنمی آپلاستیک، هپاتوتوکسیسیته.
  \item فنی‌توئین:
  \begin{itemize}
    \item مصرف یک بار در روز / \en{hL: 24–36 h}.
    \item مصرف طولانی: رویش موی اضافی، خشن‌شدن چهره، هایپرتروفی لثه.
    \item به همین دلیل در بیماران جوان نمی‌دهیم.
    \item سطح درمانی: \en{10–20 $\mu$g/mL}.
    \item در صورت مصرف تزریقی: نباید در محلول قندی انفوزیون شود \szarrow{} رسوب می‌کند.
  \end{itemize}
  \item سایر داروها: لاموتریژین، گاباپنتین، توپیرامات، تیاگابین، فنوباربیتال.
\end{itemize}

\szpage{۳}{pic03.jpg — 20261016\_205425083.jpg}
\subsection*{لاموتریژین}
\begin{itemize}
  \item معمولاً به صورت تک‌دارویی.
  \item صرع پارشیل / سندرم‌های صرعی (\en{JME}، سندرم لنوکس گاستو).
  \item عوارض: راش پوستی / تشدید میوکلونوس در \en{JME} / بالا بردن خلق و بیداری.
  \item تداخلات: اگر دارو را به والپروات اضافه می‌کنیم باید دوز دارو کمتر باشد.
\end{itemize}
\subsection*{گاباپنتین}
\begin{itemize}
  \item در بیمارانی که دچار عوارض سایر داروها هستند.
  \item به عنوان داروی کمکی / عدم تداخل دارویی با سایر داروها.
\end{itemize}
\subsection*{فنوباربیتال}
\begin{itemize}
  \item اگر در دسترس جایگزین بهتری وجود نداشته باشد.
  \item عوارض: خواب‌آلودگی در بزرگسالان، بیش‌فعالی در کودکان، تغییر شناختی.
\end{itemize}
\begin{itemize}
  \item داروی انتخابی تشنج‌های پارشیل \szarrow{} کاربامازپین.
  \item تظاهر کلاسیک تشنج پارشیل ساده \szarrow{} حرکات تونیک ـ کلونیک در یک سمت بدن.
  \begin{itemize}
    \item معمولاً صورت، گردن، اندام‌ها.
  \end{itemize}
  \item ۲۰–۳۰\% بیماران صرع به درمان طبی مقاوم هستند.
  \item شایع‌ترین اندیکاسیون جراحی: صرع لوب تمپورال.
  \item بخش انترومدیال لوب تمپورال یا به طور محدودتر هیپوکامپ و آمیگدال برداشته می‌شود.
\end{itemize}

\szpage{۴}{pic04.jpg — 20261016\_205430161.jpg}
\subsection*{تشنج ابسنس}
\begin{itemize}
  \item در ردهٔ تشنج‌های ژنرالیزهٔ اولیه.
  \item شروع در سن ۴–۸ سالگی یا اوایل نوجوانی.
  \item تظاهرات: از دست دادن ناگهانی هوشیاری برای ۵–۱۰ ثانیه.
  \begin{itemize}
    \item کنترل اندام‌ها از دست نمی‌رود.
    \item خیره‌شدن چشم‌ها.
    \item عدم وجود فاز \en{postictal}.
  \end{itemize}
  \item علائم خفیف حرکتی: پلک‌زدن سریع، جویدن، پرش کوتاه و سریع دست.
  \item تا صدها حمله در روز \szarrow{} افت تحصیلی.
  \item هایپرونتیلاسیون \szarrow{} بروز حملات.
  \begin{itemize}
    \item امواج غیرطبیعی در \en{EEG}.
  \end{itemize}
  \item \en{EEG}: امواج سوزنی جمعی دوطرفه با فرکانس ۳ \szarrow{} \en{Spike and Wave}.
\end{itemize}
\subsection*{نوع آتیپیک تشنج ابسنس}
\begin{itemize}
  \item اختلال هوشیاری طولانی‌تر.
  \item بارزتر بودن علائم حرکتی.
  \item مشخص‌نبودن زمان شروع و خاتمهٔ تشنج.
  \item متفاوت‌بودن \en{EEG}.
  \item معمولاً اختلالات ساختاری منتشر مغز وجود دارند.
  \item پاسخ بد به درمان.
\end{itemize}
\subsection*{درمان تشنج ابسنس}
\begin{itemize}
  \item اتوسوکسیماید، اسید والپروات.
  \item درمان کمکی: استازولاماید.
\end{itemize}

\szpage{۵}{pic05.jpg — 20261016\_205434052.jpg}
\subsection*{اتوسوکسیماید}
\begin{itemize}
  \item خط اول ابسنس ساده / \en{hL: 30–40 h} / سطح پلاسما: \en{50–100 $\mu$g/mL}.
  \item در درمان حرکات تونیک ـ کلونیک و موضعی بی‌اثر است.
  \item عوارض: به ندرت سرکوب \en{BM} \szarrow{} \en{CBC} دوره‌ای لازم است.
\end{itemize}
\subsection*{والپروات سدیم}
\begin{itemize}
  \item تشنج ابسنس / داروی انتخابی در بیمارانی که مخلوطی از انواع حملات صرع دارند.
  \item در موارد ابسنس + تونیک ـ کلونیک داروی انتخابی است.
\end{itemize}
\subsection*{تشنج تونیک کلونیک جنرالیزه (گرندمال)}
\begin{itemize}
  \item ناگهانی، بدون اورا و علائم هشداردهنده؛ برخی بیماران چند ساعت قبل از حمله علائم انذار (پرودرومال) دارند.
  \item شایع‌ترین علت: اختلالات متابولیک.
\end{itemize}
\subsubsection*{مراحل}
\begin{enumerate}[label=\en{\alph*)}]
  \item تونیک: عدم هوشیاری، انقباض تونیک عضلات به مدت ۱۰–۳۰ ثانیه.
  صدای فریاد یا جیغ به دلیل انقباض عضلهٔ تنفسی، سیانوز، افزایش فعالیت سمپاتیک (میدریاز، \en{RR $\uparrow$}، \en{BP $\uparrow$}).
  \item کلونیک: انقباض و انبساط متناوب عضلات / زیر ۱ دقیقه.
  \item کوما: بیمار بدون پاسخ / عضلات شل / تنفس صدادار (استریدور) به دلیل ترشح زیاد بزاق / بی‌اختیاری ادرار و مدفوع / حدود ۵ دقیقه.
  \item بهبودی: چند دقیقه تا چند ساعت / بیمار تدریجاً به هوش می‌آید.
  کنفیوژن پس از تشنج / سردرد، خستگی و درد عضلانی تا چند ساعت.
\end{enumerate}
\begin{itemize}
  \item اگر بیمار ضایعات عصبی داشته باشد \szarrow{} اختلال هوشیاری پس از تشنج ساعت‌ها طول می‌کشد.
  \item والپروات \szarrow{} خط اول درمان.
\end{itemize}

\szpage{۶}{pic06.jpg — 20261016\_205438438.jpg}
\subsection*{اندیکاسیون‌های والپروات سدیم}
\begin{itemize}
  \item ژنرالیزهٔ اولیّهٔ تونیک ـ کلونیک: خط اول.
  \begin{itemize}
    \item سپس لاموتریژین، کاربامازپین، فنی‌توئین.
  \end{itemize}
  \item حملات صرع ایدیوپاتیک.
  \item برخی تشنج‌های پارشیل (ثانویه ژنرالیزه).
  \item ابسنس، میوکلونیک، آتونیک.
  \item داروی انتخابی بیماران سندرم‌های صرعی جنرالیزه که مخلوطی از انواع حملات صرع دارند.
\end{itemize}
\subsection*{عوارض سدیم والپروات}
\begin{itemize}
  \item سمیت کبدی، سرکوب \en{BM} \szarrow{} هر دو قابل برگشت.
  \item بیماری کبدی و سابقهٔ سرکوب \en{BM} \szarrow{} کنتراندیکاسیون تجویز سدیم والپروات.
  \item نارسایی کبدی کشندهٔ غیرقابل بازگشت:
  \begin{itemize}
    \item یک عارضهٔ ایدیوسینکراتیک / در کودکان کمتر از ۲ سال.
  \end{itemize}
\end{itemize}
\subsection*{تشنج آتونیک}
% The source ends at this heading; no explanatory text follows it.
% All six sheets include the printed Date/Subject/weekday labels reproduced above.
% Handwritten page numbers 1--6 are preserved in the page headings.
\ifnum\SeizureStandalone=1 \expandafter\enddocument\fi
